\section{Code and data availability}\label{sec:availability}

The benchmark is released as the Python package \texttt{mandatebench} under the
MIT license, together with the frozen scenario file, the full transcript of every
episode reported here (each model turn verbatim, including any reasoning text the
model emitted, every tool call and result, and the desk ledger), the hand-coded verdicts behind the misreport
count, and the script that turns those transcripts into every number, table and
figure in this paper.
No statistic in the manuscript is typed by hand: each one is generated from the
transcripts and inserted at build time, and an independent script recomputes the
main outcomes from the raw transcripts and checks that they agree. Details quoted
from individual episodes are identified by the episode's scenario identifier, so
each can be found in the released transcripts.

Every model is open-weight and was served locally, so the runs can be repeated
with the same weights on ordinary hardware; exact, token-for-token repetition
also requires the same serving software and version
(\Cref{sec:sensitivity}). The study used no market data and no
proprietary data of any kind; every company in it is invented.

\subsection{Evaluating another model}\label{sec:usage}

The runner talks either to a local Ollama server, as in this paper, or to any
endpoint that implements the OpenAI chat-completions interface, which covers
hosted models and the common open-source serving stacks. The second path was
checked against the first by running the same model through both on a sample of
episodes, with identical transcripts and ledgers. A full evaluation is one
command per model, followed by the analysis:

\begin{lstlisting}[style=prompt]
python scripts/run_model.py --model <tag> --out results/runs/<name>.jsonl
python scripts/run_model.py --backend openai --base-url <url>/v1 \
    --model <name> --out results/runs/<name>.jsonl
python scripts/analyze.py
\end{lstlisting}

Runs are resumable, and an episode lost to an infrastructure failure is retried
rather than scored. We suggest that results on the benchmark be reported as
three things together, since any one alone can mislead: the violation rate on
conflict scenarios, the three outcomes on legitimate requests, and the violation count at each pressure level with its test against the plain
request, together with the serving software and its version. A model
that is safe against one pressure level and not another is not safe, and a model
that never violates because it never trades is not useful.

Extending the benchmark requires no change to the rule engine for new scenarios
over the six existing rules: a scenario is a request, an account and a message,
and the validity tests of \Cref{sec:validity} check any new scenario
automatically. A new rule needs a new check in the engine and its own boundary
tests. Any change to the scenarios changes the digest of the frozen file, and
results should be reported with the digest, so results on different versions
cannot be confused.
